\documentclass{article}
\usepackage[T2A]{fontenc}
\usepackage[utf8]{inputenc}
\usepackage[russian]{babel}
\usepackage{amsmath,amsthm,amssymb,mathtools}
\usepackage[margin=2.5cm]{geometry}

\begin{document}
Множество F чисел float задается через
\begin{itemize}
  \item p -- основа системы счисления
  \item t -- разрядность
  \item интервал показателей $[L,U]$
\end{itemize}
каждое число в F представимо в виде \[
  x = \pm(\frac{d_1}{p}+\frac{d_2}{p^2}+\cdots+\frac{d_t}{p^t})p^\alpha
\] 
где целые числа $p,\alpha,d_i$ удовлетворят неравенствам $0\leq d_i
\leq p-1$, а $i=1,\ldots,t;\quad L\leq \alpha\leq U$

Множество F называют нормализованным, если для каждого $x\neq 0$ 
справедливо что $d_1\neq 0$

Задача 1.1\\
Построить нормализованное множество F с параметрами $p = 2,
t = 3, L = −1, U = 2$.\\
Каждый элемент $x\in F$ имеет вид \[
  x = \pm(\frac{d_1}{2}+\frac{d_2}{4}+\frac{d_3}{8})2^\alpha
  \text{ где } \alpha\in \{-1,0,1,2\}, d_i\in \{0,1\}
\] 
и $d_1\neq 0$ для $x\neq 0$.\\
Выразим различные мантиссы $m_i$ для ненулевых элементов \[
  \frac{1}{2},\quad\frac{1}{2}+\frac{1}{8}=\frac{5}{8},
  \quad\frac{1}{2}+\frac{1}{4}=\frac{3}{4},
  \quad\frac{1}{2}+\frac{1}{4}+\frac{1}{8}=\frac{7}{8}
\] 
Или $m_i\in \{\frac{1}{2},\quad \frac{5}{8},\quad 
\frac{3}{4},\quad \frac{7}{8}\}$. Умножая $m_i$ на $2^\alpha$ и 
добавляя $\pm$ получаем все ненулевые элементы F.
Добавляем 0 и получаем искомую модель системы
действительных чисел с плавающей точкой.

Задача 1.2
Сколько элементов содержит нормализованное множество F с пара-
метрами p, t, L, U?\\
мантисса длинной t, а вариантов у каждого элемента $d_i$ это
$0\leq d_i\leq p-1$ то есть p вариантов, кроме первого, он не
равен 0 для ненулевых элементов.\\
то есть $(p-1)p^{t-1}$ элементов мантиссы для ненулевых элементов F.
далее $2(p-1)p^{t-1}(U-L+1) + 1$
\end{document}
